\subsection{Theorem~\ref{thm:cache_topology} Proof: Optimizing \Sys cache topology}

To prove Theorem~\ref{thm:cache_topology}, we will first prove Theorem~\ref{thm:cache_topology_general}, which is a more general version of the theorem.
Then, Theorem~\ref{thm:cache_topology} will be an easy corollary of this more general theorem.


\begin{theorem}
\label{thm:cache_topology_general}
The choice of $F_k^p$ (and equivalently, $F_k^b$) values that maximizes the speedup of \Sys under the constraint $\sum_{k=0}^K F_k^p \leq B$ (where $K$ is the speculative lookahead), is:
\begin{eqnarray*}
\sum_{k=0}^K F_k^p &=& B, \\
\frac{\partial p_{\mathrm{hit,p}}^{k}}{\partial F}\big(F_k\big) &=& a_p^{-k}\cdot \frac{\partial p_{\mathrm{hit,p}}^{0}}{\partial F}\big(F_0\big) \quad \forall k < K, \;\; \text{and} \\
\frac{\partial p_{\mathrm{hit,p}}^{K,all}}{\partial F}\big(F_K\big) &=& (1-a_p) \cdot a_p^{-K}\cdot \frac{\partial p_{\mathrm{hit,p}}^{0}}{\partial F}\big(F_0\big). \\
% p_l'(F_K) &=& (1-a_p) a^{-K}\frac{\partial }{\partial F_k} p_{\mathrm{hit,p}}(F_0)
\end{eqnarray*}
\end{theorem}


\begin{proof}


To understand how to optimize our cache topology, we first rewrite the speedup equations from Section~\ref{sec:ssd_theory}, now making explicit how the speedup depends on the $F_k^p$ and $F_k^b$ values:
\begin{eqnarray*}
speedup\Big( \{F_k^p\}_{k=0}^K), \{F_k^b\}_{k=0}^K) \Big) &=& p_{\mathrm{hit}} \Big(\{F_k^p\}, \{F_k^b\}\Big) \cdot E_{\mathrm{hit}} \\
 &\;& + \; \; \;  \bigg(1-p_{\mathrm{hit}} \Big(\{F_k^p\}, \{F_k^b\}\Big)\bigg) \cdot E_{\mathrm{miss}}, \quad \text{where} \\
p_{\mathrm{hit}} \Big(\{F_k^p\}, \{F_k^b\}\Big) &=& \frac{p_{\mathrm{hit,b}}(\{F_k^b\})}{1 + p_{\mathrm{hit,b}}(\{F_k^b\}) - p_{\mathrm{hit,p}}(\{F_k^p\}) }.
\end{eqnarray*}
We can express $p_{\mathrm{hit,p}}$ more precisely in terms of:
\begin{itemize}
    \item The acceptance rate $a_p$ of the primary speculator, and
    \item The functions $p_{\mathrm{hit,p}}^{k}(F_k)$ and $p_{\mathrm{hit,p}}^{K,all}(F_k)$:
    These functions describe the probability of a cache hit conditioned on (1) the last round's speculator was the primary one, (2) $k$ tokens were accepted, (3) whether all of the tokens were accepted ($p_{\mathrm{hit,p}}^{K,all}$) or not ($p_{\mathrm{hit,p}}^{k}$), and (4) $F_k$ verification outcomes were prepared for in case $k$ tokens were accepted.
    Note that this can be estimated empirically by comparing the draft and target probability distributions for a set of sequences in a calibration set.
\end{itemize}
We can do the same for $p_{\mathrm{hit,b}}$, the cache hit rate when the last round's speculator was the backup.
\begin{eqnarray*}
p_{\mathrm{hit,p}}(\{F_k^p\}) &=& a_p^K  \cdot p_{\mathrm{hit,p}}^{K,all}(F_K) + \sum_{k=0}^{K-1} a_p^k (1-a_p) \cdot p_{\mathrm{hit,p}}^{k}(F_k), \quad {and} \\
% p_{\mathrm{hit,b}}(\{F_k^b\}) &=& a_b^K  \cdot p_{\mathrm{hit,b}}^{K,all}(F_K) + \sum_{k=0}^{K-1} a_b^k (1-a_b) \cdot p_{\mathrm{hit,b}}^{k}(F_k).
\end{eqnarray*}

These equations are a direct result of the fact that the chance of accepting exactly $k$ tokens is $a_p^k(1-a_p)$ for $k < K$, and $a_p^K$ for $k=K$, when the acceptance rate is $a_p$.

% By Lemma~\ref{lemma:phit_simplify}

% We first prove that $p_{\mathrm{hit,p}}(F)$ and $p_{\mathrm{hit,b}}(F)$ are equal to the following:
% \begin{eqnarray*}
% p_{\mathrm{hit,p}}(F) &=& a_p^K  \cdot p_{\mathrm{hit,p}}^{K,all}(F_K) + \sum_{k=0}^{K-1} a_p^k (1-a_p) \cdot p_{\mathrm{hit,p}}(F_k) \\
% &=& p_{\mathrm{hit,p}}(F) + a_p^K \cdot \Big( p_{\mathrm{hit,p}}^{K,all}(F) -  p_{\mathrm{hit,p}}(F) \Big)\quad \text{if we assume $F_k = F$ for all $k$.} \\
% p_{\mathrm{hit,b}}(F) &=& a_b^K  \cdot p_{\mathrm{hit,b}}^{K,all}(F_K) + \sum_{k=0}^{K-1} a_b^k (1-a_b) \cdot p_{\mathrm{hit,b}}(F_k) \\
% &=& p_{\mathrm{hit,b}}(F) + a_b^K \cdot \Big( p_{\mathrm{hit,b}}^{K,all}(F) -  p_{\mathrm{hit,b}}(F) \Big)\quad \text{if we assume $F_k = F$ for all $k$.} \\
% \end{eqnarray*}

% \noindent The first equality (in both cases) is true because:
% \begin{itemize}
% \item The chance of accepting exactly $k$ tokens is equal to
% $a_p^K (1-a_p)$ for $0 \leq k < K$, and equal to $a_p^K$ when $k=K$.
% (where $a$ equals $a_p$ for a cache hit, and $a_b$ for a cache miss), and 
% \item Whenever we accept exactly $k$ tokens, the probability of getting a cache hit is $p_{\mathrm{hit,n/nn}}(F_k)$ (for $k<K$) and $p_{\mathrm{hit,n/nn}}(F_k)^{K,all}$ by definition.
% \end{itemize}

% \noindent The second equality (in both cases) is a result of a telescoping sum, where all middle terms cancel out. \\
% \begin{eqnarray*}
% \sum_{k=0}^{K-1} a_p^K (1-a_p) &=& 1 - a + a - a^2 + a^2 - a^3 \ldots - a^{K-1} + a^{K-1} - a_p^K \\
% &=& 1 - a_p^K
% \end{eqnarray*}

% \noindent We can then rearrange to get the desired expressions:
% \begin{eqnarray*}
% a_p^K p_{\mathrm{hit,p}}^{K,all}(F) + (1-a_p^K) \cdot p_{\mathrm{hit,p}}(F) &=& p_{\mathrm{hit,p}}(F) + a_p^K \cdot \Big(p_{\mathrm{hit,p}}^{K,all}(F) - p_{\mathrm{hit,p}}(F) \Big) \\
% a_p^K p_{\mathrm{hit,b}}^{K,all}(F) + (1-a_p^K) \cdot p_{\mathrm{hit,b}}(F) &=& p_{\mathrm{hit,b}}(F) + a_p^K \cdot \Big(p_{\mathrm{hit,b}}^{K,all}(F) - p_{\mathrm{hit,b}}(F) \Big)
% \end{eqnarray*}

Given these equations, we can now optimize the cache topology (i.e., the choice of $F_k$ values) to maximize speedup.
We notice that the speedup is monotonically increasing in $p_{\mathrm{hit}}$, and that $p_{\mathrm{hit}}$ is monotonically increasing in both $p_{\mathrm{hit,p}}$ and $p_{\mathrm{hit,b}}$ (easy to see by taking first derivatives of $p_{\mathrm{hit}}$, and seeing they are always positive).
Thus, we must simply maximize $p_{\mathrm{hit,p}}(\{F_k^p\})$ and $p_{\mathrm{hit,b}}(\{F_k^b\})$ under the constraints that $\sum_{k=0}^K F_k^p \leq B$ and $\sum_{k=0}^K F_k^b \leq B$.
We do this below.

\subsubsection{Maximizing $p_{\mathrm{hit,p}}(\{F_k^p\})$ and $p_{\mathrm{hit,b}}(\{F_k^b\})$}

\noindent As discussed, we want to maximize the probability of a cache hit under the budget constraint $\sum_{k=0}^K F_k \leq B$.
We do this for $p_{\mathrm{hit,p}}$, and the proof is identical for $p_{\mathrm{hit,b}}$.
\begin{eqnarray*}
% \max_{\sum_k F_k \leq B} p_{\mathrm{hit,p}}(K,F) = \max_{\sum_k F_k \leq B} \sum_{k=0}^K a_p^k (1-a_p)^{\ind{k < K}} \cdot p_{\mathrm{hit,p}}(F_k) 
\max_{\sum_k F_k \leq B} p_{\mathrm{hit,p}}(K,F) = \max_{\sum_k F_k \leq B} a_p^K  \cdot p_{\mathrm{hit,p}}^{K,all}(F_K) + \sum_{k=0}^{K-1} a_p^k (1-a_p) \cdot p_{\mathrm{hit,p}}^k(F_k)
\end{eqnarray*}

\noindent  We will solve this maximization problem with Lagrange multipliers.
\begin{eqnarray*}
\mathcal{L}(F_0, \ldots, F_k, \lambda) = a_p^K  \cdot p_{\mathrm{hit,p}}^{K,all}(F_K) + \sum_{k=0}^{K-1} a_p^k (1-a_p) \cdot p_{\mathrm{hit,p}}^k(F_k) + \lambda \cdot \Big(\sum_{k=0}^K F_k - B\Big) \\
% \Rightarrow \frac{\partial \mathcal{L}}{\partial F_k} = a_p^k(1-a_p)^{\ind{k < K}} \cdot p_{\mathrm{hit,p}}'(F_k) = 0
% \frac{\partial }{\partial F_k} p_{\mathrm{hit,p}}(F_k) = 0 
\end{eqnarray*}

\noindent Now, we will take the derivative with respect to all the variables, and set it to zero.

\begin{eqnarray*}
\frac{\partial \mathcal{L}}{\partial F_k} &=& a_p^k(1-a_p) \cdot \frac{\partial }{\partial F_k} p_{\mathrm{hit,p}}^k(F_k) + \lambda = 0  \quad \text{for } k < K \\
\frac{\partial \mathcal{L}}{\partial F_K} &=& a_p^K \frac{\partial }{\partial F_k} p_{\mathrm{hit,p}}^{K,all}(F_K) + \lambda = 0 \\
\frac{\partial \mathcal{L}}{\partial \lambda} &=& \sum_{k=0}^K F_k - B = 0 \\
% \frac{\partial }{\partial F_k} p_{\mathrm{hit,p}}(F_k) = 0 
\end{eqnarray*}

\noindent We notice that for all $k$, $\frac{\partial \mathcal{L}}{\partial F_k}$ are equal to one another (all equal to $-\lambda$). Thus, we see that:
\begin{eqnarray*}
(1-a_p)\frac{\partial }{\partial F_k} p_{\mathrm{hit,p}}^0(F_0) &=& a_p(1-a_p)\frac{\partial }{\partial F_k} p_{\mathrm{hit,p}}^1(F_1) = \ldots \\
\ldots &=&  a_p^{K-1}(1-a_p)\frac{\partial }{\partial F_k} p_{\mathrm{hit,p}}^{K-1}(F_{K-1}) = a_p^K \frac{\partial }{\partial F_k} p_{\mathrm{hit,p}}^{K,all}(F_K) \\
\Rightarrow \frac{\partial }{\partial F_k} p_{\mathrm{hit,p}}^k(F_k) &=& a_p^{-k}\frac{\partial }{\partial F_k} p_{\mathrm{hit,p}}^0(F_0) \quad \forall k < K, \;\; \text{and} \\
\frac{\partial }{\partial F_k} p_{\mathrm{hit,p}}^{K,all}(F_K) &=& (1-a_p) a_p^{-K}\frac{\partial }{\partial F_k} p_{\mathrm{hit,p}}^0(F_0)
\end{eqnarray*}

This gives the desired result.

% Intuitively, this theorem shows that because the chance of accepting exactly $k$ tokens is greater than the chance of accepting $k' > k$ tokens, we should use more of our budget on $F_k$ than on $F_{k'}$.

% A practical challenge in applying this theorem is that we do not have closed-form expressions for $p_{\mathrm{hit,p}}^{k}$ and $p_{\mathrm{hit,p}}^{K,all}$, and can only measure them empirically at integer values of $F_k$.
% Thus, we must estimate the derivatives of these functions (e.g., using finite-difference formulas, or by fitting cubic splines) in order to solve the above equations.
% \AM{
% Let's lower bound the expected speedup, as a function of $B$ and $b$.
% \begin{enumerate}
% \item Assume cache miss rate EQUALS $b$ power law $\Rightarrow$ pick optimal $F_k$
% \item Plug those into $p_{\mathrm{hit}}$ equation to lower bound $p_{\mathrm{hit}}$ as a function of $b$, $B$ (here, we will assume  $p_{\mathrm{hit}}$ has a $r$ power-law cache hit rate (cache miss rate LOWER-BOUNDED by power law).
% \item Assume $p_{\mathrm{hit,b}} \geq p_{\mathrm{hit,p}}$, then the expected speedup is $\geq p_{\mathrm{hit}} \cdot E_{\mathrm{hit}}$.
% \end{enumerate}
% }

% \TK{
% The bound on expected speedup is great, but should probably go in the Appendix. 
% }

\end{proof}

We now use this result to prove (an extended version of) Theorem~\ref{thm:cache_topology}, which is a special case of the above theorem in the case where the speculators have $r$ power-law cache hit rates.

% \begin{theorem}
\paragraph{Theorem~\ref{thm:cache_topology}. (\textit{\Sys Cache Shape: Geometric Fan-Out})}
\label{thm:cache_topology_app}
% \textbf{(\Sys Cache Topology)}
\textit{The optimal choice of $F_k^p$ (equivalently, $F_k^b$) values for $k\in [0,K]$ for \Sys under the constraint $\sum_{k=0}^K F_k^p \leq B$, and under the assumption that the speculator has acceptance rate $a_p$ and a $r$ power-law cache hit rate, follows a geometric series (for $k < K$):}
\begin{eqnarray*}
F_k &=& F_0 \cdot a_p^{k/(1 + r)}  \;\; \forall k < K, \; \textit{and} \\
F_K &=& F_0 \cdot a_p^{K/(1 + r)} \cdot (1-a_p)^{-1/(1 + r)},
\end{eqnarray*}
\text{where $F_0$ can be chosen as follows so that $\sum_{k=0}^K F_k^p = B$.}
\begin{eqnarray*}
    F_0 &=& \frac{B}{a_p^{K/(1 + r)} \cdot (1-a_p)^{-1/(1 + r)} +  \Big(1-a_p^{K/(1 + r)}\Big)/\Big( 1-a_p^{1/(1 + r)}\Big)} 
\end{eqnarray*}
An equivalent result holds for $F_k^b$.
% \end{theorem}

% \begin{proof}
% For simplicity, we'll just let $a \coloneqq a_p$, $p_{\mathrm{hit,p}}F_k) \coloneqq p_{\mathrm{hit,p}}(F_k)$, $\frac{\partial }{\partial F_k} p_{\mathrm{hit,p}}(F_k) \coloneqq \frac{\partial }{\partial F_k} p_{\mathrm{hit,p}}(F_k)$, and $p_l'(F_k) \coloneqq \frac{\partial }{\partial F_k} p_{\mathrm{hit,p}}^{K,all}(F_k)$.
% For simplicity, we'll just let $a \coloneqq a_p$, $p_k(F_k) \coloneqq p_{\mathrm{hit,n}}^k(F_k)$, $p_k'(F_k) \coloneqq \frac{\partial }{\partial F_k} p_{\mathrm{hit,n}}^k(F_k)$, and $p_{K,all}'(F_K) \coloneqq \frac{\partial }{\partial F_K} p_{\mathrm{hit,n}}^{K,all}(F_K)$.

\textit{Proof.} Now, substituting $p_{\mathrm{hit,p}}^k(F) = 1 - F^{-r}$ (and thus $\frac{\partial }{\partial F_k} p_{\mathrm{hit,p}}^k(F) = r \cdot F^{-r-1}$),  we get:
\begin{eqnarray*}
\Rightarrow r \cdot F_k^{-r-1} &=& a_p^{-k}r \cdot F_0^{-r-1} \;\; \forall k < K, \; \text{and} \\
r \cdot F_K^{-r-1} &=& (1-a_p) \cdot a_p^{-K} \cdot r \cdot F_0^{-r-1}. \\
\Rightarrow F_k &=& F_0 \cdot a_p^{k/(1 + r)}  \;\; \forall k < K, \; \text{and} \\
F_K &=& F_0^{(1+r)/(1+r)} \cdot a_p^{K/(1 + r)} \cdot (1-a_p)^{-1/(1 + r)} \cdot (r/r)^{-1/(1+r)}.
\end{eqnarray*}

To solve for the exact sequence of $F_k$ fan-out values, we plug the above values into the budget equation:
\begin{eqnarray*}
    B &=& F_0 \cdot a_p^{K/(1 + r)} \cdot (1-a_p)^{-1/(1 + r)}  + \sum_{k=0}^{K-1} F_0 \cdot a_p^{k/(1 + r)} \\
\end{eqnarray*}

Solving for $F_0$ gives:
\begin{eqnarray*}
    F_0 &=& \frac{B}{ a_p^{K/(1 + r)} \cdot (1-a_p)^{-1/(1 + r)} + \sum_{k=0}^{K-1} a_p^{k/(1 + r)}} \\
\end{eqnarray*}
We can simplify this further using the equation for the sum of the geometric series:
$$\sum_{k=0}^{K-1} a_p^{k/(1 + r)} = \sum_{k=0}^{K-1} c^k = \frac{1-c^K}{1-c}, \quad \text{for } c = a_p^{1/(1 + r)}$$

Plugging this in gives the desired result.
\begin{eqnarray*}
    F_0 &=& \frac{B}{a_p^{K/(1 + r)} \cdot (1-a_p)^{-1/(1 + r)} +  \Big(1-a_p^{K/(1 + r)}\Big)/\Big( 1-a_p^{1/(1 + r)}\Big)} 
\end{eqnarray*}
% \end{proof}

$\hfill\square$